\section{Discussion}

The study distinguishes three properties that are often conflated. Verdict correctness asks whether the final label matches the benchmark. Explanation quality asks whether the prose is useful and readable. Evidence attribution asks whether the cited evidence supports and materially influences the verdict. A system may satisfy any two while failing the third.

Counterfactual tests strengthen an attribution claim but do not reveal a model's complete internal causal process. Deletion can create out-of-distribution inputs, and an unchanged prediction may reflect redundant valid evidence. The protocol therefore deletes complete admissible evidence groups, reports sufficiency alongside comprehensiveness, and analyzes naturally redundant cases separately.

The public-data design trades broad social coverage for reproducibility and expert evidence annotations. General-domain transfer to MOCHEG tests whether the approach is confined to scientific documents, but it cannot establish global or multilingual generalization.
